
\section{10. Limitations}

\subsection{10.1 Task Coverage}

Only SST-2 evaluated ($L = 128$). ASA's sparsity benefit grows with $L$; 
we do not measure accuracy at $L \ge 512$.

\subsection{10.2 Model Coverage}

Tested on models with $H = 12$. Not tested on:
\begin{itemize}
\item Very wide models ($H = 32+$)
\item Mixture-of-Experts
\item Vision transformers
\end{itemize}

\subsection{10.3 No Parameter Reduction}

ASA \textbf{does not} reduce parameter count. It reduces attention memory and 
compute. For parameter savings, use A3 or AASVD.

\subsection{10.4 No Time Reduction}

Chunking bounds memory but does not reduce \textbf{time} -- we still compute all 
scores. To reduce time, we would need to skip score computation for pruned 
keys.

\subsection{10.5 Approximate Causality}

Causal masking is applied to each chunk. The chunk boundary may leak 
information from prior chunks (in bidirectional settings), but for causal 
decoders, we verify strict causality.

\subsection{10.6 Sparsity Logit Floor}

The floor operation in $k_h = \lfloor (1 - \rho_h) L \rfloor$ is not 
differentiable. We use a straight-through estimator, which works but is 
not theoretically justified.

\subsection{10.7 Single Seed}

All results from seed 42.

\subsection{10.8 No Comparison with Some Baselines}

We do not directly compare with:
\begin{itemize}
\item Longformer
\item BigBird
\item Routing Transformer
\item Reformer
\end{itemize}

\subsection{10.9 Chunk Boundary Effects}

For very small chunks ($C = 32$), chunk boundary effects may hurt accuracy.

\subsection{10.10 Failure Modes}

Causality, long-context evaluation

\subsection{10.11 Code Quality}

Research-grade implementation.

\subsection{10.12 Ethical Considerations}

Dual-use. See Appendix E.